\documentclass[border=5pt]{standalone}

% --- Paquetes de Codificación y Lenguaje ---
\usepackage[utf8]{inputenc}
\usepackage[T1]{fontenc}
\usepackage[spanish, es-tabla]{babel}

% --- Paquetes Matemáticos y de Ingeniería ---
\usepackage{amsmath}
\usepackage{amsfonts}
\usepackage{amssymb}
\usepackage{siunitx} % Para magnitudes físicas

% --- Configuración de Colores ---
\usepackage{xcolor}
\definecolor{black}{HTML}{000000}

% --- Núcleo de Diagramación: TikZ y CircuiTikZ ---
\usepackage{tikz}
\usepackage[american, cute inductors, american currents, american voltages]{circuitikz}

% --- Librerías de TikZ Especializadas ---
\usetikzlibrary{arrows.meta}    % Para puntas de flecha IEEE
\usetikzlibrary{positioning}   % Para posicionamiento relativo de nodos
\usetikzlibrary{calc}          % Para cálculos de coordenadas
\usetikzlibrary{shapes.geometric}
\usetikzlibrary{backgrounds}   % PARA CAPAS DE FONDO/FOREGROUND

% --- Configuración de Estilo CircuiTikZ (Estándar IEEE) ---
\ctikzset{
	logic ports=ieee,                     % Compuertas según IEEE Std 91/91a
	resistors/scale=0.8,                  % Escalamiento profesional
	capacitors/scale=0.7,
	inductors/scale=0.7,
	diodes/scale=0.7,
	voltage/distance from node=0.4,       % Ajuste de etiquetas de voltaje
	bipoles/thickness=1                   % Grosor de línea para componentes
}

% --- Definición de Estilos Personalizados ---
\tikzset{
	label style/.style={font=\small\sffamily},
	connection/.style={circle, fill=black, inner sep=1.5pt} % Para nodos de conexión (junctions)
}

% --- MACRO OSCILOSCOPIO MATEMÁTICO (adaptada para capa background) ---
\def\oscope#1#2#3#4#5{
	\begin{scope}[shift={(#1,#2)}]
		% Chasis del osciloscopio
		\draw[fill=white, thick, rounded corners=2pt] (-0.2, -1.2) rectangle (3.2, 1.4);
		% Rejilla de la pantalla (Graticule)
		\draw[step=0.25, gray!30, thin] (0, -1) grid (3, 1);
		% Ejes X e Y
		\draw[gray!80, thin] (0, 0) -- (3, 0); 
		\draw[gray!80, thin] (0, -1) -- (0, 1);
		% Etiqueta del instrumento
		\node[above, font=\sffamily\tiny\bfseries] at (1.5, 1.05) {#3};
		% Trazado de la onda (path matemático puro)
		\begin{scope}[shift={(0,0)}]
			\draw[thick, black] #5;
		\end{scope}
		% Punto de anclaje de la sonda (solo coordenada interna)
		\coordinate (osc_bot) at (1.5, -1.2);
	\end{scope}
}

\begin{document}
	
	\begin{circuitikz}[line width=0.8pt, american]
		
		% =======================================================
		% CAPA DE FONDO (BACKGROUND): 5 OSCILOSCOPIOS MATEMÁTICOS
		% =======================================================
		\begin{pgfonlayer}{background}
			
			% 1. Entrada Red AC (senoidal completa)
			\oscope{-2.5}{9.5}{1. Entrada Red AC}{N1}{
				plot[domain=0:3.2, samples=120] (\x, {0.85*sin(2.5*\x r)})
			}
			
			% 2. Secundario del transformador (senoidal reducida)
			\oscope{1.8}{9.5}{2. Secundario Transformador}{N2}{
				plot[domain=0:3.2, samples=120] (\x, {0.42*sin(2.5*\x r)})
			}
			
			% 3. Rectificado de onda completa
			\oscope{6.1}{9.5}{3. Rectificado Puente}{N3}{
				plot[domain=0:3.2, samples=180] (\x, {0.42*abs(sin(2.5*\x r))})
			}
			
			% 4. Rizado + Offset (descarga del bulk)
			\oscope{10.4}{9.5}{4. Rizado + Offset DC}{N4}{
				(0,0.28) -- (0.55,0.42) -- (1.05,0.24) -- (1.65,0.39) -- (2.15,0.23) -- (2.75,0.37) -- (3.2,0.28)
			}
			
			% 5. Salida DC regulada (línea plana)
			\oscope{14.7}{9.5}{5. DC Regulada Estable}{N5}{
				(0,0.18) -- (3.2,0.18)
			}
			
		\end{pgfonlayer}
		
		% =======================================================
		% CAPA FRONTAL (FOREGROUND): ESQUEMÁTICO UNIFICADO 3 NIVELES
		% =======================================================
		
		% --- NIVEL 1 (MACRO): Etapa de red AC, transformador, rectificador, filtrado y carga ---
		% Transformador
		\node[transformer core] (T) at (0,0) {};
		\draw (T.A1) to[short, -o] ++(-1.8,0) coordinate (N1) node[left] {$V_{ac,+}$};
		\draw (T.A2) to[short, -o] ++(-1.8,0) node[left] {$V_{ac,-}$};
		
		% Puente rectificador discreto
		\coordinate (Bleft) at (2.5, 0);
		\coordinate (Bright) at (4.5, 0);
		\coordinate (Btop) at (3.5, 1.5);
		\coordinate (Bbot) at (3.5, -1.5);
		
		\draw (Bleft) to[D] (Btop);
		\draw (Bright) to[D] (Btop);
		\draw (Bbot) to[D] (Bleft);
		\draw (Bbot) to[D] (Bright);
		
		% Nodos físicos del puente
		\node[circle, fill, inner sep=1.2pt] at (Bleft) {};
		\node[circle, fill, inner sep=1.2pt] at (Bright) {};
		\node[circle, fill, inner sep=1.2pt] (N3) at (Btop) {};
		\node[circle, fill, inner sep=1.2pt] at (Bbot) {};
		
		% Conexiones transformador → puente
		\draw (T.B1) -| (Bleft);
		\draw (T.B2) -- ++(0.5,0) |- (4.5, -2.5) -- (Bright);
		
		% Rieles de alimentación principales
		\draw (Btop) -- (3.5, 2.5) -- (7.5, 2.5) coordinate (RailIn); % IN del regulador Nivel 2
		\draw (Bbot) -- (3.5, -3.5) -- (19.5, -3.5);
		\draw (19, -3.5) node[ground] {};
		
		% Filtrado previo (bulk + cerámico)
		\draw (5.5, 2.5) node[circle, fill, inner sep=1.2pt] (N4) {} to[eC, l=$C_{bulk}$, *-*] (5.5, -3.5);
		\draw (6.7, 2.5) to[C, l=$C_{in}$, *-*] (6.7, -3.5);
		
		% --- NIVEL 2 (REGULADOR A NIVEL SILICIO): Integrado en el riel positivo (sin caja) ---
		% IN = RailIn (x=7.5, y=2.5)
		
		% Riel de entrada extendido (Vcc interno)
		\draw (7.5, 2.5) -- (8.5, 2.5) coordinate (INnode);
		
		% Transistor de paso (Q_pass NPN) en serie
		\node[npn] (Qpass) at (10.5, 4.2) {};
		\node[above left=2pt] at (Qpass.B) {$Q_{pass}$};
		
		% Colector conectado al riel de entrada
		\draw (INnode) -- (Qpass.C);
		
		% Resistencia de sensado de corriente (R_sc) en emisor
		\draw (Qpass.E) node[circle, fill, inner sep=1.2pt] (NE) {} to[R, l=$R_{sc}$, *-] (12.8, 2.5) coordinate (SenseOut);
		
		% Salida del regulador (OUT)
		\draw (SenseOut) -- (16.5, 2.5) coordinate (N5);
		
		% Etapa de protección contra sobrecorriente (Q_limit NPN)
		\node[npn] (Qlim) at (14.2, 5.1) {};
		\node[right=2pt] at (Qlim.B) {$Q_{limit}$};
		
		% Sensado diferencial sobre R_sc
		\draw (NE) -- ++(0,1.1) coordinate (SenseTop) -- (14.2, 5.1) -- (Qlim.B);
		\draw (Qlim.E) -- (13.5, 2.5);
		
		% El colector de Q_limit roba corriente de la base de Q_pass
		\draw (Qlim.C) -- ++(0,0.8) -| (Qpass.B) node[circle, fill, inner sep=1.2pt] (NB) {};
		
		% Referencia de voltaje interna (bias + Zener)
		\draw (8.2, 2.5) to[R, l=$R_{bias}$] (8.2, 0.8) to[zzD, l=$V_{ref}$, -*] (8.2, -1.2) coordinate (VrefNode);
		\draw (VrefNode) -- (8.2, -3.5);
		
		% Red de retroalimentación (divisor R1/R2) desde OUT
		\draw (16.5, 2.5) -- (16.5, 0.5) to[R, l=$R_1$, *-*] (16.5, -0.8) coordinate (FBtap) to[R, l=$R_2$, -*] (16.5, -3.5);
		
		% --- NIVEL 3 (AMPLIFICADOR DE ERROR A NIVEL TRANSISTOR): Integrado dentro del Nivel 2 ---
		% Par diferencial de entrada NPN (Q1/Q2)
		\node[npn] (Q1) at (9.1, 1.1) {}; \node[below right=1pt] at (Q1.B) {$Q_1$};
		\node[npn, xscale=-1] (Q2) at (11.1, 1.1) {}; \node[below left=1pt] at (Q2.B) {$Q_2$};
		
		% Entradas: + → Vref, - → FBtap
		\draw (VrefNode) -- ++(0,0.6) -| (Q1.B);
		\draw (FBtap) -- ++(0.4,0) |- (Q2.B);
		
		% Espejo de corriente de carga activa PNP (Q3/Q4)
		\node[pnp] (Q3) at (9.1, 2.8) {}; \node[above right=1pt] at (Q3.B) {$Q_3$};
		\node[pnp, xscale=-1] (Q4) at (11.1, 2.8) {}; \node[above left=1pt] at (Q4.B) {$Q_4$};
		
		% Emisores a Vcc (INnode)
		\draw (Q3.E) -- (9.1, 2.8) -- (INnode |- 9.1, 2.8) node[circle, fill, inner sep=1.2pt] {};
		\draw (Q4.E) -- (11.1, 2.8) -- (INnode |- 11.1, 2.8) node[circle, fill, inner sep=1.2pt] {};
		
		% Colectores y nodo diferencial (V_diff)
		\draw (Q1.C) -- (Q3.C);
		\draw (Q2.C) -- (Q4.C) node[midway, circle, fill, inner sep=1.2pt] (Vdiff) {};
		
		% Conexión bases espejo (diodo Q3)
		\draw (Q3.B) -- (8.4, 2.8) -- (8.4, 3.4) -- (11.8, 3.4) -- (11.8, 2.8) -- (Q4.B);
		\draw (8.4, 2.8) |- (Q1.C) node[circle, fill, inner sep=1.2pt] {};
		
		% Fuente de corriente de cola (I_tail)
		\draw (Q1.E) -- (Q2.E) -- (10.1, 0.4) node[circle, fill, inner sep=1.2pt] {} to[I, l=$I_{tail}$] (10.1, -1.2) node[circle, fill, inner sep=1.2pt] {};
		
		% Etapa VAS (Q5 NPN)
		\node[npn] (Q5) at (12.8, 0.4) {}; \node[below right=1pt] at (Q5.B) {$Q_5$};
		\draw (Vdiff) -- (12.1, 2.1) |- (Q5.B);
		\draw (Q5.E) -- (12.8, -1.1) node[circle, fill, inner sep=1.2pt] {};
		
		% Polarización y diodos clase AB (D1/D2)
		\draw (INnode) -- (13.5, 3.8) to[I, l=$I_{bias}$] (13.5, 2.9) node[circle, fill, inner sep=1.2pt] (NodeA) {};
		\draw (NodeA) to[D, l=$D_1$] (13.5, 1.8) to[D, l=$D_2$] (13.5, 0.8) node[circle, fill, inner sep=1.2pt] (NodeB) {};
		\draw (NodeB) -- (Q5.C);
		
		% Capacitor de compensación Miller (Cc)
		\draw (12.1, 0.4) node[circle, fill, inner sep=1.2pt]{} -- (12.1, 1.6) to[C, l=$C_c$, -*] (13.5, 1.6);
		
		% Etapa de salida Push-Pull clase AB (Q6 NPN / Q7 PNP)
		\node[npn] (Q6) at (15.1, 3.8) {}; \node[above left=1pt] at (Q6.B) {$Q_6$};
		\node[pnp] (Q7) at (15.1, 0.8) {}; \node[below left=1pt] at (Q7.B) {$Q_7$};
		
		% Alimentación colectores salida
		\draw (Q6.C) -- (15.1, 4.5) -- (INnode |- 15.1, 4.5) node[circle, fill, inner sep=1.2pt] {};
		\draw (Q7.C) -- (15.1, -1.1) -- (VrefNode |- 15.1, -1.1) node[circle, fill, inner sep=1.2pt] {};
		
		% Conexiones bases Push-Pull
		\draw (NodeA) -- (Q6.B);
		\draw (NodeB) -- (Q7.B);
		
		% Resistencias de degeneración de emisor y nodo OUT del error amp
		\draw (Q6.E) to[R, l=$R_{e1}$] (15.1, 2.4) node[circle, fill, inner sep=1.2pt] (ERRout) {};
		\draw (Q7.E) to[R, l=$R_{e2}$] (ERRout);
		
		% Salida del error amp → nodo de drive (NB) del Q_pass
		\draw (ERRout) -- (14.1, 4.0) -| (NB);
		
		% --- FILTRADO DE SALIDA Y CARGA (continuación Nivel 1) ---
		\draw (16.8, 2.5) to[C, l=$C_{out}$, *-*] (16.8, -3.5);
		\draw (17.8, 2.5) to[eC, l=$C_{load}$, *-*] (17.8, -3.5);
		\draw (18.8, 2.5) to[R, l=$R_{load}$, *-*] (18.8, -3.5);
		
		% Puertos de salida finales
		\draw (19, 2.5) to[short, -o] (19.8, 2.5) node[right] {$V_{out,+}$};
		\draw (19, -3.5) to[short, -o] (19.8, -3.5) node[right] {$V_{out,-}$};
		
		% Conexión de tierra del regulador (Nivel 2/3)
		\draw (8.2, -1.2) -- (16.5, -3.5);
		
		% Diodo de protección flyback (OUT → IN)
		\draw (16.5, 2.5) ++(0.3,0) coordinate (P_OUT) to[short, *-] ++(0,1.4) to[D, l=$D_{prot}$] ++(-7.5,0) to[short, -*] (INnode);
		
	\end{circuitikz}
	
\end{document}